% Artigo da Vitrine — template ACM sigconf (Overleaf: "ACM Conference Proceedings").
% Compilar no Overleaf: subir a pasta artigo/ inteira e usar pdfLaTeX + BibTeX.
\documentclass[sigconf,nonacm]{acmart}

\usepackage[brazilian]{babel}
\usepackage{listings}
\usepackage{booktabs}

\lstdefinelanguage{vitrine}{
  morekeywords={produto,regra,quando,entao,carrinho,calcular,desconto,frete,em,qtd,total,x},
  morecomment=[l]{\#},
  sensitive=true
}
\lstset{language=vitrine, basicstyle=\ttfamily\footnotesize, keywordstyle=\bfseries,
        columns=fullflexible, frame=single, inputencoding=utf8,
        literate={ã}{{\~a}}1 {á}{{\'a}}1 {é}{{\'e}}1 {ç}{{\c{c}}}1 {õ}{{\~o}}1 {ê}{{\^e}}1}

\title{Vitrine: regras de promoção como linguagem compilada}

\author{João Pedro Gaspar}
\author{Juliard Martins}
\author{Nathan Kevyn}
\author{Lucas Nogueira}
\author{Kauê Pablo}
\affiliation{%
  \institution{Centro Universitário Adventista de São Paulo (UNASP)}
  \city{São Paulo}
  \country{Brasil}}

\begin{abstract}
% TODO (todos, por último): 4–6 frases. Problema, o que a Vitrine faz,
% a verificação central, o principal número da avaliação, a principal limitação.
\end{abstract}

\keywords{linguagens de domínio específico, análise estática, compiladores, comércio eletrônico}

\begin{document}
\maketitle

% DONO: todos (fechar juntos no fim)   ·   ESPAÇO: ~0,75 página
\section{Introdução}
% O enunciado pede: o domínio, o problema e por que uma linguagem própria
% faz sentido EM VEZ DE UMA BIBLIOTECA.
%
% Pontos a cobrir:
% - Promoções de loja virtual são regras de negócio que se combinam.
%   Combinações não previstas (descontos que se somam, fretes que se
%   contradizem) só aparecem quando um carrinho real as dispara.
% - Plataformas validam um carrinho por vez, em execução. CUIDADO: não
%   afirmem nada sobre Shopify/VTEX sem citar a documentação delas
%   (ex.: regras de combinação de descontos). Não digam que o total fica
%   negativo: as plataformas limitam o preço a zero.
% - Tese: com uma linguagem de condições FECHADA (só qtd, total e
%   constantes), o compilador analisa o conjunto de regras para QUALQUER
%   carrinho possível. Numa biblioteca Python, a condição seria uma função
%   arbitrária, e não dá para decidir se duas funções arbitrárias valem
%   juntas. Citar Mernik et al. \cite{mernik2005} sobre quando criar uma DSL.
% - Contribuições (lista curta): a linguagem; a verificação de conflitos
%   por faixas; a avaliação com N programas.
% - Uma frase sobre a organização do artigo.

% DONO: João Pedro (exemplo, construções) + Nathan (decisões)   ·   ESPAÇO: ~1,25 página
\section{A linguagem}
% O enunciado pede: programa de exemplo, as construções principais e TRÊS
% decisões de projeto justificadas, cada uma com a alternativa descartada.

\subsection{Exemplo}
% Usar exemplos/validos/promocao_outubro.vitrine e o recibo. Explicar em
% prosa: produto, regra quando/entao, carrinho, calcular.
\begin{lstlisting}
produto camiseta = R$ 49,90;
produto caneca   = R$ 29,90;
frete = R$ 15,00;
regra leve3 { quando qtd(camiseta) >= 3
              entao desconto 20% em camiseta; }
regra frete_gratis { quando total >= R$ 150,00
                     entao frete R$ 0,00; }
carrinho { camiseta x 3; caneca x 1; }
calcular;
\end{lstlisting}

\subsection{Construções e tipos}
% - Tipos: inteiro, dinheiro (centavos inteiros), percentual (só em desconto).
% - Precedência pela cadeia expressao > termo > fator. Ex.:
%   R$ 10,00 + R$ 2,00 * 3 = R$ 16,00.
% - Escopo global único; declarar antes de usar.
% - Uma tabela pequena de tipos (docs/especificacao.md, seção 5).

\subsection{Decisões de projeto}
% Três decisões (guia, seção 6). Para cada uma: decisão, alternativa, por quê.
% 1. Descontos aditivos sobre o preço original x descontos compostos.
%    SEJAM HONESTOS: com composição a soma nunca passa de 100%, então a
%    escolha cria o invariante que a verificação protege. Digam por que
%    ainda assim faz sentido (previsibilidade para o lojista).
% 2. total = subtotal antes dos descontos x total corrente na ordem das
%    regras (que tornaria o resultado dependente da ordem e invalidaria a
%    análise do conjunto). Caso concreto: no exemplo canônico, R$ 179,60
%    antes do desconto e R$ 149,66 depois, com frete grátis a partir de
%    R$ 150,00.
% 3. Análise conservadora por faixas x soma ingênua (recusa faixas
%    exclusivas, o caso mais comum) x solver exato (SMT: biblioteca pronta,
%    proibida pelo enunciado para a semântica).

\section{Implementação}
% DONO: João Pedro   ·   ESPAÇO: ~0,25 página
\subsection{Léxico e sintático}
% - Lark \cite{lark}, parser LALR, a partir de grammar.lark; Transformer
%   para a AST em dataclasses.
% - OBRIGATÓRIO pelo enunciado: o que se ganhou e o que se perdeu com o
%   gerador. Ganhou: gramática legível igual à EBNF, linha e coluna prontas.
%   Perdeu: controle das mensagens de erro (a lista de tokens esperados é
%   imperfeita em alguns estados LALR).
% - Uma dificuldade real encontrada por VOCÊ (ex.: palavra-chave x
%   identificador, prioridade de PERCENTUAL sobre NUMERO, e.accepts x
%   e.expected) e como foi resolvida. Se fizer a tarefa do separador de
%   milhar, contar aqui.

% DONO: Juliard   ·   ESPAÇO: ~0,25 página
\subsection{Análise semântica}
% - Tabela de símbolos com escopo global; "declarado antes" por posição.
% - Verificação de tipos recursiva; preço como expressão constante avaliada
%   na compilação (R$ 1,00 / 0 é erro semântico, não de execução).
% - Para no primeiro erro (por quê) e, se fizer a tarefa, a versão que
%   relata todos os erros de referência.
% - Uma dificuldade real e como foi resolvida.

% DONO: Nathan   ·   ESPAÇO: ~0,5 página (é a contribuição central)
\subsection{Verificação de conflitos}
% - Cada condição "qtd(p) op c" ou "total op c" vira uma faixa de inteiros
%   (domínio de intervalos; citar Cousot e Cousot \cite{cousot1977}).
% - Viabilidade: interseção das faixas + acoplamento
%   total = soma de qtd(p) * preço(p). Exemplo: qtd(camiseta) >= 3 implica
%   total >= R$ 149,70, incompatível com total < R$ 100,00.
% - Pior combinação: busca com/sem cada regra e duas podas (a inviabilidade
%   só cresce com o conjunto; limite pela soma restante). Exponencial no
%   pior caso: medir o tempo com 20, 40 e 80 regras.
% - Condições não analisáveis viram "pode valer sempre" (conservador): nunca
%   aceita um conflito real, mas pode recusar um programa correto.
% - Se fizer as tarefas (c04, m02): o que mudou no algoritmo.

% DONO: Lucas   ·   ESPAÇO: ~0,25 página
\subsection{Execução}
% - Dinheiro em centavos inteiros, nunca float; arredondamento meio para
%   cima por linha do carrinho (ex.: 15% de 3 x R$ 49,90 = 2245,5 -> 2246
%   centavos; por unidade daria 2247).
% - Independência da ordem das regras (há um teste que inverte a ordem).
% - Testes: N testes pytest, testes de recibo esperado e CI no GitHub Actions.
% - Compilador no navegador: o MESMO código Python via Pyodide
%   \cite{pyodide}, sem reimplementação (uma frase só).

% DONO: Kauê (+ números do Nathan)   ·   ESPAÇO: ~1,5 página — A SEÇÃO DE MAIOR PESO
\section{Avaliação}
% Opção B do enunciado: a verificação semântica funciona?
%
% Método:
% - N programas rotulados em avaliacao/casos/ (meta: 20 a 30), cada um com
%   gabarito "aceita" ou "recusa <categoria>".
% - Separar SEMENTES (13 casos gerados com IA, que já incluíam falhas
%   conhecidas) dos CASOS DA EQUIPE (escritos à mão, sem olhar o código).
%   Isso precisa estar explícito no texto.
% - Métricas: detectados, escapes, falsos positivos, aceitos corretamente;
%   tempo médio de verificação.
%
% Resultados: copiar a tabela gerada por `python avaliacao/rodar.py`.
% Resultado atual só das sementes (27/10), como referência:
%   detecção 6/7 (86%), falsos positivos 2/6 (33%), 1 escape, ~0,65 ms.
%
% Interpretação (o slide 20 exige):
% - De onde vêm os falsos positivos (c04: forma não normalizada; c05:
%   integralidade ignorada) e o escape (m02).
% - Antes x depois das melhorias do Nathan.
% - O que os números NÃO permitem concluir: os casos não vêm de lojas reais,
%   a amostra é pequena, os autores escreveram os casos e o verificador.
\begin{table}[h]
  \caption{Resultados da verificação semântica (preencher)}
  \begin{tabular}{lrrrr}
    \toprule
    categoria & casos & detectado & escape & falso pos. \\
    \midrule
    conflito\_desconto & & & & \\
    conflito\_frete & & & & \\
    regra\_morta & & & & \\
    controles (aceita) & & & & \\
    \bottomrule
  \end{tabular}
\end{table}

% DONO: Juliard   ·   ESPAÇO: ~0,5 página
\section{Limitações}
% O enunciado: o que não faz ou faz de forma frágil, COM A CAUSA de cada uma.
% Dizer que não há limitações dá teto de 4,0 neste critério.
% Base: docs/especificacao.md, seção 8. Escolher as mais importantes:
% - falsos positivos por formas não analisáveis e por integralidade;
% - uma condição por regra, sem e/ou;
% - o invariante de 100% é fraco (plataformas já limitam a zero); um piso de
%   margem seria mais útil;
% - busca exponencial no número de regras por produto;
% - para no primeiro erro; tokens mínimos (sem separador de milhar etc.).

% DONO: todos   ·   ESPAÇO: um parágrafo
\section{Conclusão}
% O que foi construído, se funciona e como sabemos (o número principal), e
% o próximo passo natural (ex.: piso de margem; gerar a configuração de uma
% plataforma de e-commerce como alvo de compilação).

% DONO: todos   ·   OBRIGATÓRIO (omitir zera o critério de integridade)
\section*{Declaração de uso de inteligência artificial}
% O enunciado pede: ferramenta, em quais etapas e com que grau de supervisão.
% Rascunho factual (ajustar ao que for verdade no fim):
% - Ferramenta: Claude (Anthropic).
% - Etapas: geração do código inicial do compilador (léxico, sintático,
%   semântico, execução), dos testes, dos exemplos, das 13 sementes da
%   avaliação e do compilador web; revisão da especificação; esqueleto deste
%   artigo (estrutura e tópicos de cada seção). Um protótipo anterior em
%   JavaScript também foi gerado com IA.
% - Supervisão: a equipe revisou, executou e modificou o código (citar as
%   mudanças de cada integrante); os casos 14 em diante da avaliação e o
%   texto do artigo foram escritos pela equipe.


% \nocite{*} lista todas as entradas do .bib mesmo sem \cite no texto.
% Apagar esta linha quando as seções já citarem as referências de verdade.
\nocite{*}
\bibliographystyle{ACM-Reference-Format}
\bibliography{referencias}

\appendix
% DONO: Lucas
\section{Reprodução}
% - Repositório: https://github.com/joaoplgaspar/vitrine
% - Terminal: python -m venv .venv; pip install -e ".[dev]"; pytest;
%   python -m vitrine exemplos/validos/promocao_outubro.vitrine;
%   python avaliacao/rodar.py
% - Navegador: link do GitHub Pages (sem instalação).


\end{document}
